\documentclass[12pt,a4paper]{article}

\usepackage{fontawesome5}
\providecommand{\faPencilAlt}{\faIcon{pencil-alt}}

\usepackage{styles/couleurs}
\usepackage{styles/mise_en_page}
\usepackage{styles/boites}
\usepackage{styles/maths}
\usepackage{styles/figures}
\usepackage{styles/evaluations}

\usepackage{enumitem}
\usepackage{array}
\usepackage{tabularx}
\usepackage{amsmath}
\usepackage{amssymb}

\setlist[itemize]{leftmargin=1.4em,itemsep=0.15em,topsep=0.2em}
\setlist[enumerate]{leftmargin=*,itemsep=0.25em,topsep=0.25em}

\begin{document}

\enteteinterro[Calculatrice interdite]{4e}{Facteurs premiers et fractions}{18 min}{12}

\exo{4}[Calculer]
Décompose en produit de facteurs premiers. \pts{4}

\[
84\qquad 126\qquad 150
\]

\lignes{6}

\exo{4}[Calculer]
Simplifie à l'aide des facteurs premiers. \pts{4}

\[
\frac{42}{63}
\qquad
\frac{84}{126}
\]

\lignes{6}

\exo{2}[Raisonner]
Parmi les fractions suivantes, lesquelles sont égales ?
Justifie. \pts{2}

\[
\frac{6}{15}
\qquad
\frac{10}{25}
\qquad
\frac{15}{36}
\]

\lignes{3}

\exo{2}[Communiquer]
Une élève écrit :

\[
90=2\times3\times15.
\]

Explique pourquoi ce n'est pas encore une décomposition
en produit de facteurs premiers, puis termine-la. \pts{2}

\lignes{3}

\end{document}
